\subsection{官方 DGPPO 的任务适配与预算控制比较}\label{subsec:external}
为增加独立的学习方法参照，采用 Zhang 等\cite{zhang2025dgppo}公开的 DGPPO 实现，保持其图 Transformer 策略、任务价值网络、分解约束价值网络、离散障碍优势和近端策略更新不变。沿用上游提供的无循环网络选项，在独立环境接口中将二阶运动改为本文单积分器；约束接口分别输入碰撞、全局谱连通和工作区的有符号代价。训练时保留原实现的代价裕度变换；该变换不应用于物理距离、状态积分或评测审计。

两种方法均使用 $N\in\{6,\ldots,10\}$ 的原始格点任务生成规则。DGPPO 从该生成器建立每规模 $4096$ 个任务的有限库，按生成器接受样本的规模频率抽取同规模小批次；EGP-NPA 使用原有在线任务采样。DGPPO 每个展开窗口为 $500$ 步，每个训练重复共 $480000$ 个环境转移，其中随机采样和更新内部的确定性展开均计入预算。两个 EGP-NPA 训练重复分别使用 $481990$ 与 $484821$ 个环境步。两者均在测试前指定最终预算检查点，不按测试结果选模。

外部算法保留自身参数化和安全学习目标：DGPPO 使用 $6$ 维节点、$5$ 维边输入，含相对目标、相对几何、边界余量及精确谱裕度；EGP-NPA 使用 $17/20$ 维输入并包含自身护盾的内部状态。DGPPO 的任务奖励采用归一化距离进展和时间步代价，未加入 EGP-NPA 的对偶相关停滞奖励及到达终止奖励。网络结构、内部状态与学习目标的差异属于完整控制流程的组成，分别披露，而不把该比较解释为仅替换一个学习更新公式的因果实验。

DGPPO 策略、任务价值网络、约束价值网络分别含 $37796$、$33441$、$14211$ 个参数；EGP-NPA 策略含 $108998$ 个参数。本轮保持原生网络规模，未人为匹配参数量；因此性能差异同时包含表示、控制结构及训练配置的影响。

测试采用另外生成并冻结的 $96$ 个任务：$N=8$ 为 $64$ 个，$N=12,20,25$ 分别为 $12,12,8$ 个。两个训练重复、两种方法共执行 $384$ 次闭环，使用相同初始状态、目标、安全距离、速度上限、时间上限、到达保持条件和独立有理数审计。DGPPO 直接执行策略速度，不叠加本文护盾；EGP-NPA 的分配与保持段认证属于其系统组成，认证及预算终止仍单独记录。

\begin{table*}[t]
\centering\footnotesize
\caption{新生成96任务上的外部方法适配比较。平均完成时间将未安全到达计为60秒；审计失败不自动等同于物理碰撞。}\label{tab:external}
\begin{tabularx}{\textwidth}{@{}Xccccc@{}}\toprule
方法 & 重复 & $N=8$ 安全到达 & 跨规模安全到达 & 审计失败运行 & 平均完成时间/s\\\midrule
EGP-NPA & 1 & 64/64 & 29/32 & 0/96 & 13.67 \\
EGP-NPA & 2 & 64/64 & 28/32 & 0/96 & 13.67 \\
DGPPO（任务适配） & 1 & 0/64 & 0/32 & 96/96 & 60.00 \\
DGPPO（任务适配） & 2 & 0/64 & 0/32 & 96/96 & 60.00 \\
\bottomrule\end{tabularx}
\end{table*}

EGP-NPA第1次训练的名义到达为$93/96$，安全到达为$93/96$，发生物理碰撞的运行有$0$次。
EGP-NPA第2次训练的名义到达为$92/96$，安全到达为$92/96$，发生物理碰撞的运行有$0$次。
DGPPO第1次训练的名义到达为$0/96$，安全到达为$0/96$，发生物理碰撞的运行有$96$次。
DGPPO第2次训练的名义到达为$0/96$，安全到达为$0/96$，发生物理碰撞的运行有$96$次。


表\ref{tab:external}评价上述实现和预算下的端到端结果。当前每种方法仅有两个训练重复，有限任务库与在线采样、训练奖励以及既有调参投入未完全匹配。DGPPO 的全局谱约束属于本任务扩展，其表现不能替代原论文环境中的结论。本轮不据此报告算法级显著性或宣称一般性优越；若适配策略未满足安全条件，应将其解释为当前训练配置未达到任务要求，而非所有离散障碍学习方法必然失败。
